%%%%%%%%%%%%%%%%%%%%%%%%%%%%%%%%%%%%%%%%%%%%%%%%%%%%%%%%%%%%
\section{结论}
\label{sec:conclusion}

我们提出了认证工作流（authenticated workflows）——面向企业级智能体 AI 的首个完整信任层。我们的关键架构洞察是：定位在协议层——介于在组合场景下失效的应用层防御与缺乏工作流语义的基础设施原语之间——能够在对异构框架提供统一保护的同时，保持工作流上下文。所有智能体与资源的交互均经过四个基本控制面（提示、工具、数据、上下文）；统一保护这些控制面，实现了在任一相邻层单独无法达成的组合安全性。

三项贡献实现了这一愿景：MAPL 通过密码学证言提供密码学工作流依赖，从而在无需信任应用层报告状态的前提下实现顺序约束；嵌入式独立 PEP（策略执行点）实现纵深防御，攻击者必须同时攻破多个密码学屏障；通过薄适配器（adapter）集成九个框架，验证了控制面的完备性。形式化证明（引理 1--7、定理 1--3）建立了安全属性；实证验证在 174 个测试用例上展示了 100\% 召回率、零假阳性，以及对两个攻破基线系统的生产环境 CVE 的防护。这解决了第~\ref{sec:introduction} 节提出的五项挑战：跨框架的密码学完整性、规模化策略执行、组合过程中的权限不提升、通过哈希链（hash chain）实现的上下文完整性、以及通过不可否认性（non-repudiation）实现的完全可问责性。随着智能体 AI 向生产环境过渡，认证工作流提供了安全基底，使得在组合缺口当前阻碍采用的场景下能够安全部署。
